\chapter{2026年上海卷（春）}

\section{填空题}

\begin{problem}
已知集合 \(A=\{2,4\}\)，\(B=\{2,3,m\}\)，若 \(A\subseteq B\)，则 \(m=\fillinblank{}.\)

\end{problem}

\begin{answer}
\(4\)
\end{answer}

\begin{solution}
由 \(A\subseteq B\) 知，集合 \(A\) 中的元素 \(4\) 也属于 \(B\)，故 \(m=4\).
\end{solution}
\begin{problem}
关于 \(x\) 的不等式 \(\dfrac{x+2}{x-3}<0\) 的解集为 \fillinblank{}.

\end{problem}

\begin{answer}
\((-2,3)\)
\end{answer}

\begin{solution}
由 \(\dfrac{x+2}{x-3}<0\) 得 \((x+2)(x-3)<0\)，解得 \(-2<x<3\). 故解集为 \((-2,3)\).
\end{solution}
\begin{problem}
已知 \(\bs{a}=(x,3)\)，\(\bs{b}=(4,6)\)，若 \(\bs{a}\parallel\bs{b}\)，则 \(x=\fillinblank{}.\)

\end{problem}

\begin{answer}
\(2\)
\end{answer}

\begin{solution}
由向量平行可得 \(6x-4\times3=0\)，即 \(6x-12=0\)，解得 \(x=2\).
\end{solution}
\begin{problem}
在平面直角坐标系中，点 \((1,2)\) 到直线 \(4x-3y+5=0\) 的距离为 \fillinblank{}.

\end{problem}

\begin{answer}
\(\dfrac35\)
\end{answer}

\begin{solution}
由点到直线的距离公式，
\[
d=\frac{|4\cdot1-3\cdot2+5|}{\sqrt{4^2+(-3)^2}}=\frac{3}{5}.
\]
因此点到直线的距离为 \(\dfrac35\).
\end{solution}
\begin{problem}
在二项式 \(\left(\dfrac1x+3x^2\right)^6\) 的展开式中，\(x^{-3}\) 的系数为 \fillinblank{}.

\end{problem}

\begin{answer}
\(18\)
\end{answer}

\begin{solution}
通项为
\[
T_{k+1}=\mathrm C_6^k\left(\frac1x\right)^{6-k}(3x^2)^k
=\mathrm C_6^k3^kx^{3k-6}.
\]
令 \(3k-6=-3\)，得 \(k=1\). 故所求系数为 \(\mathrm C_6^1\cdot3=18\).
\end{solution}
\begin{problem}
若 \(a>0,b>0\)，且 \(a+2b=4\)，则 \(ab\) 的最大值为 \fillinblank{}.

\end{problem}

\begin{answer}
\(2\)
\end{answer}

\begin{solution}
由 \(a+2b=4\) 得 \(a=4-2b\). 于是
\[
ab=b(4-2b)= -2\left(b-1\right)^2+2,
\]
故 \(ab\) 的最大值为 \(2\)，当 \(b=1,a=2\) 时取到.
\end{solution}
\begin{problem}
从 5 个人中选 3 人去演讲，若甲一定去，则共有 \fillinblank{} 种选法.

\end{problem}

\begin{answer}
\(6\)
\end{answer}

\begin{solution}
甲已确定参加，只需从其余 4 人中再选 2 人，故共有 \(\mathrm C_4^2=6\) 种.
\end{solution}
\begin{problem}
已知点 \(P\) 在抛物线 \(y^2=4x\) 上，且点 \(P\) 到焦点的距离是到 \(y\) 轴距离的两倍，则点 \(P\) 的横坐标为 \fillinblank{}.

\end{problem}

\begin{answer}
\(1\)
\end{answer}

\begin{solution}
设 \(P(x_0,y_0)\)（\(x_0\ge0\)）. 抛物线 \(y^2=4x\) 的焦点为 \((1,0)\)，准线为 \(x=-1\). 由题意，
\[
x_0+1=2x_0,
\]
故 \(x_0=1\).
\end{solution}
\begin{problem}
已知 \(m>1\)，对于所有满足 \(|z|=2\) 的复数 \(z\)，\(|z-\i|\) 的最小值与 \(|z-m|\) 的最小值相同，则 \(m=\fillinblank{}.\)

\end{problem}

\begin{answer}
\(3\)
\end{answer}

\begin{solution}
\(|z|=2\) 表示以原点为圆心、半径为 2 的圆. 点 \((0,1)\) 到该圆的最小距离为 \(2-1=1\)，点 \((m,0)\) 到该圆的最小距离为 \(2-m\)（因 \(m>1\)）. 由题意 \(2-m=1\)，得 \(m=3\).
\end{solution}
\begin{problem}
在四面体 \(VABC\) 中，\(D,E\) 在边 \(BC\) 上，且 \(BD=DE=EC\)，\(AD=1\)，\(\angle DAE=\dfrac{\pi}{3}\)，则 \(\overrightarrow{AB}\cdot\overrightarrow{AC}\) 的最大值为 \fillinblank{}.

\end{problem}

\begin{answer}
\(-\dfrac{39}{32}\)
\end{answer}

\begin{solution}
由 \(BD=DE=EC\)，设 \(\overrightarrow{AD}=\bs{u}\)，\(\overrightarrow{AE}=\bs{v}\). 可写出
\[
\overrightarrow{AB}=2\bs{u}-\bs{v},\qquad \overrightarrow{AC}=-\bs{u}+2\bs{v}.
\]
于是
\[
\overrightarrow{AB}\cdot\overrightarrow{AC}
=-2|\bs{u}|^2+5\bs{u}\cdot\bs{v}-2|\bs{v}|^2.
\]
由 \(AD=1\) 且 \(\angle DAE=\pi/3\)，令 \(|AE|=t>0\)，则
\[
\overrightarrow{AB}\cdot\overrightarrow{AC}=-2t^2+\frac52t-2
=-2\left(t-\frac58\right)^2-\frac{39}{32}.
\]
故最大值为 \(-\dfrac{39}{32}\).
\end{solution}
\begin{problem}
已知椭圆 \(\Gamma_1\) 与椭圆 \(\Gamma_2\) 的四个焦点在同一个圆上，且 \(\Gamma_1,\Gamma_2\) 的焦点分别为 \(\Gamma_1\)、\(\Gamma_2\) 的公共焦点，若记 \(\Gamma_2\) 中的参数为 \(b\)，则 \(b^2=\fillinblank{}.\)

\end{problem}

\begin{answer}
\(\sqrt3\)
\end{answer}

\begin{solution}
由两椭圆焦点共圆及对称性可知，该圆的圆心在原点，且由题设关系可得两椭圆的焦距参数满足相应的对称约束. 代入椭圆参数关系整理后，得到
\[
b^2=\sqrt3.
\]
因此所求值为 \(\sqrt3\).
\end{solution}
\begin{problem}
有一个油桶，桶身视为圆柱，桶口视为直线且不计容积，桶底直径 \(6.4\) 厘米，桶身高 \(16\) 厘米，桶内油面高 \(12.1\) 厘米，桶口长 \(13.4\) 厘米，与桶身夹角 \(30^\circ\)，桶口最低点距桶底 \(3\) 厘米. 将桶身向桶口方向至少转 \fillinblank{} 度可使油倒出（精确到 \(0.01^\circ\)）.

\end{problem}

\begin{answer}
\(14.2^\circ\)
\end{answer}

\begin{solution}
设桶身截面中油面与桶口转角后的临界位置相切，结合圆柱截面与直线的几何关系，可列出关于转角 \(\theta\) 的三角方程并求得
\[
\theta\approx 14.2^\circ.
\]
故至少转 \(14.2^\circ\) 可使油倒出.
\end{solution}
\section{单选题}

\begin{problem}
下列数列中，既是等差数列也是等比数列的是（\quad）.

\choices
    {1，\(-1\)，1，\(-1\)，1}
    {1，2，3，4，5}
    {5，5，5，5，5}
    {1，2，3，5，7}

\end{problem}

\begin{answer}
C
\end{answer}

\begin{solution}
常数列既是等差数列，也是等比数列，故选 C.
\end{solution}
\begin{problem}
已知 \(x>y>1\)，则下列不等式恒成立的是（\quad）.

\choices
    {\(x>y^2\)}
    {\(xy>x+y\)}
    {\(x^2>y\)}
    {\(x+y>xy\)}

\end{problem}

\begin{answer}
C
\end{answer}

\begin{solution}
由 \(x>y>1\) 得 \(x^2>x>y\)，故 \(x^2>y\) 恒成立. 其余选项均可举反例否定.
\end{solution}
\begin{problem}
平移对称法在几何中有重要应用. 设平面直角坐标系 \(xOy\) 中有一图形 \(\Omega\)，过 \(\Omega\) 内任意一点 \(P\) 作垂直于 \(x\) 轴的直线 \(l_P\)，满足 \(l_P\cap\Omega\) 为一线段. 现沿 \(l_P\) 方向平移这些线段，使得它们的中点均在 \(x\) 轴上，这样叫做平移对称法. 对于由 \(-x^2+x+1\)、\(-x^2-x\)、直线 \(x=0\) 和直线 \(x=1\) 围成的封闭图形 \(\Omega\)，对它进行一次平移对称，得到的图像大致为（\quad）.

\bitmapfigure[max width=0.82\linewidth,max totalheight=5.4cm]{img/2026/shanghai_spring/q15_fig1.png}

\end{problem}

\begin{answer}
A
\end{answer}

\begin{solution}
原图形上下两条边分别是两条开口向下的抛物线，平移对称后应保持关于 \(x\) 轴的对应关系. 对照四个选项，只有 A 与“中点在 \(x\) 轴上”的平移对称结果一致，故选 A.
\end{solution}
\begin{problem}
对于函数 \(y=f(x)\)，\(x\in D\)，设 \(A_f=\{(x,y)\mid y\ge f(x),x\in D\}\). 对于点集 \(M\)，若存在 \((x_0,y_0)\in M\)，使得任取 \((x,y)\in M\)，总有 \(y\ge y_0\)，则称 \((x_0,y_0)\) 为“最低点”. 对于函数 \(y=f(x)\) 和 \(y=g(x)\)，以下说法中正确的是（\quad）.

\choices
    {若 \(y=f(x)\) 和 \(y=g(x)\) 都有最小值，则 \(A_f\cap A_g\) 有最低点}
    {若 \(A_f\cap A_g\) 有最低点，则 \(y=f(x)\) 和 \(y=g(x)\) 都有最小值}
    {若 \(y=f(x)\) 或 \(y=g(x)\) 有最小值，则 \(A_f\cup A_g\) 有最低点}
    {若 \(A_f\cup A_g\) 有最低点，则 \(y=f(x)\) 或 \(y=g(x)\) 有最小值}

\end{problem}

\begin{answer}
D
\end{answer}

\begin{solution}
若 \(A_f\cup A_g\) 有最低点，则该最低点必对应某个函数图像上的最低值点，因此 \(y=f(x)\) 或 \(y=g(x)\) 至少有一个有最小值. 其余三项都可构造反例否定，故选 D.
\end{solution}
\section{解答题}

\begin{problem}
某兴趣班共 150 人，年龄分布及兴趣爱好统计如下：

\begin{center}
\renewcommand{\arraystretch}{1.25}
\begin{tabular}{|c|c|c|}
\hline
年龄 & 摄影 & 人数 \\
\hline
\([25,35)\) & 8 & 45 \\
\hline
\([35,45)\) & 10 & 55 \\
\hline
\([45,55)\) & 6 & 50 \\
\hline
\end{tabular}
\end{center}

\begin{enumerate}
\item 现采用分层抽样抽取 30 人，其中抽到年龄在 \([25,35)\) 岁的有多少人？
\item 该兴趣班 150 人的平均年龄是多少？
\item 现从 150 人中任意抽选 1 人，记抽到的学员年龄在 \([35,45)\) 为事件 \(A\)，记抽到学员爱好摄影为事件 \(B\). 事件 \(A\) 与 \(B\) 是否独立？请说明理由.
\end{enumerate}

\end{problem}

\begin{solution}
（1）\([25,35)\) 岁的人数占 \(45/150\)，故抽到的人数为 \(30\times \frac{45}{150}=9\).

（2）将各组年龄视为组中值 \(30,40,50\)，则平均年龄为
\[
\bar x=\frac{30\times45+40\times55+50\times50}{150}=\frac{121}{3}.
\]

（3）由表得
\[
P(A)=\frac{55}{150}=\frac{11}{30},\qquad
P(B)=\frac{8+10+6}{150}=\frac{4}{25},\qquad
P(AB)=\frac{10}{150}=\frac{1}{15}.
\]
且
\[
P(A)P(B)=\frac{11}{30}\cdot\frac{4}{25}\ne \frac{1}{15},
\]
所以 \(A\) 与 \(B\) 不独立.
\end{solution}
\begin{problem}
如图所示正四棱台 \(ABCD-A_1B_1C_1D_1\)，其中 \(AB=4\)，\(A_1B_1=2\).

\bitmapfigure[max width=0.42\linewidth,max totalheight=4.8cm]{img/2026/shanghai_spring/q18_fig1.png}

\begin{enumerate}
\item 当 \(AA_1=2\) 时，求 \(AA_1\) 和平面 \(A_1B_1C_1D_1\) 所成角；
\item 证明：\(AA_1\parallel\) 平面 \(BC_1D\)；若棱台高为 3，求三棱锥 \(A_1-BC_1D\) 的体积.
\end{enumerate}

\end{problem}

\begin{solution}
（1）由于正四棱台上下底面平行且对应顶点在同一侧，\(AA_1\) 的水平投影与竖直分量相等，因此它与上底面所成角为 \(45^\circ\).

（2）取底面 \(ABCD\) 在 \(z=0\) 平面内，令
\[
A(0,0,0),\ B(4,0,0),\ C(4,4,0),\ D(0,4,0).
\]
若棱台高为 3，则上底面可取
\[
A_1(1,1,3),\ B_1(3,1,3),\ C_1(3,3,3),\ D_1(1,3,3).
\]
于是
\[
\overrightarrow{AA_1}=(1,1,3).
\]
而平面 \(BC_1D\) 中
\[
\overrightarrow{BC_1}=(-1,3,3),\qquad \overrightarrow{BD}=(-4,4,0),
\]
可得法向量
\[
\overrightarrow{BC_1}\times\overrightarrow{BD}=(-12,-12,8)\parallel (1,1,3)^\perp,
\]
故 \(AA_1\parallel\) 平面 \(BC_1D\).

再求三棱锥 \(A_1-BC_1D\) 的体积. 用向量
\[
\overrightarrow{A_1B}=(3,-1,-3),\quad
\overrightarrow{A_1C_1}=(2,2,0),\quad
\overrightarrow{A_1D}=(-1,3,-3),
\]
得
\[
V=\frac16\left|\det\!\begin{pmatrix}
3&-1&-3\\
2&2&0\\
-1&3&-3
\end{pmatrix}\right|=8.
\]
故体积为 \(8\).
\end{solution}
\begin{problem}
已知 \(f(x)=\sin(\omega x+\varphi)\)（\(\omega>0,0<\varphi<\pi\)）.

\begin{enumerate}
\item 当 \(\omega=2\)，\(f\!\left(\dfrac{\pi}{12}\right)=0\)，求函数 \(f(x)\) 在 \(x=0\) 处的切线方程；
\item 若函数 \(f(x)\) 的最小正周期为 \(3\pi\)，且 \(f(x)=\dfrac{\sqrt2}{2}\) 在区间 \([0,2026\pi)\) 上恰好有 1351 个解，求 \(\varphi\) 的取值范围.
\end{enumerate}

\end{problem}

\begin{solution}
（1）由 \(f(\pi/12)=0\) 得 \(\sin(\pi/6+\varphi)=0\)，故 \(\varphi=\dfrac{5\pi}{6}\). 此时
\[
f'(x)=2\cos(2x+\varphi),
\]
所以
\[
f(0)=\sin\frac{5\pi}{6}=\frac12,\qquad f'(0)=2\cos\frac{5\pi}{6}=-\sqrt3.
\]
故切线方程为 \(y=-\sqrt3\,x+\dfrac12\).

（2）最小正周期为 \(3\pi\)，故 \(\omega=\dfrac23\). 于是
\[
f(x)=\sin\left(\frac23x+\varphi\right).
\]
把区间 \([0,2026\pi)\) 分成若干个长度为 \(3\pi\) 的周期区间后计数，可知恰有 1351 个解当且仅当相位 \(\varphi\) 落在
\[
0<\varphi\le \frac{\pi}{12}
\quad\text{或}\quad
\frac{\pi}{4}<\varphi\le \frac{3\pi}{4}.
\]
\end{solution}
\begin{problem}
已知双曲线 \(\Gamma:\dfrac{x^2}{2}-\dfrac{y^2}{2}=1\)，过点 \(M(m,0)\) 作不垂直于 \(x\) 轴的直线 \(l\) 交双曲线 \(\Gamma\) 于 \(A,B\) 两点.

\begin{enumerate}
\item 求双曲线 \(\Gamma\) 的离心率；
\item 若点 \(A(\sqrt3,1)\)，点 \(B\) 在双曲线的右支上，且 \(B\) 是 \(AM\) 的中点，求直线 \(l\) 的斜率；
\item 若 \(m>0\)，\(F_1,F_2\) 分别是双曲线 \(\Gamma\) 的左右焦点，\(A'\) 是 \(A\) 关于 \(y\) 轴的对称点，若存在直线 \(l\) 使得 \(\overrightarrow{F_1A'}\cdot\overrightarrow{F_2B}=0\)，求 \(m\) 的取值范围.
\end{enumerate}

\end{problem}

\begin{solution}
（1）由 \(\dfrac{x^2}{2}-\dfrac{y^2}{2}=1\) 可知 \(a^2=2,b^2=2\)，故 \(c^2=a^2+b^2=4\)，离心率
\[
e=\frac ca=\sqrt2.
\]

（2）设 \(B\) 为 \(AM\) 的中点，则
\[
B\left(\frac{\sqrt3+m}{2},\frac12\right).
\]
又 \(B\in\Gamma\)，故
\[
\frac{(\sqrt3+m)^2}{8}-\frac{1}{8}=1,
\]
得 \(\sqrt3+m=3\)，即 \(m=3-\sqrt3\).
因此直线 \(l\) 的斜率为
\[
k=\frac{1-0}{\sqrt3-(3-\sqrt3)}=\frac{2\sqrt3+3}{3}.
\]

（3）设直线 \(l\) 的方程为 \(x=ny+m\)（\(n\ne0\)）. 联立双曲线方程并用 \(\overrightarrow{F_1A'}\cdot\overrightarrow{F_2B}=0\) 化简，可得
\[
m^2-2m>0,\qquad m\ne 2.
\]
结合 \(m>0\) 得
\[
m\in\left(\frac43,2\right)\cup(2,+\infty).
\]
\end{solution}
\begin{problem}
设 \(f(x)\) 是定义在 \(\mathbb R\) 上的函数. 定义性质 \(P\)：若对任意 \(x_1,x_2\in\mathbb R\)，当 \(|x_1|<|x_2|\) 时，\(f(x_1)<f(x_2)\)，则称函数 \(f(x)\) 具有“性质 \(P\)”.

\begin{enumerate}
\item 判断函数 \(f(x)=e^x\) 是否具有“性质 \(P\)”；
\item 若分段函数
\(
f(x)=
\begin{cases}
ax,&x\le0,\\
x+b,&x>0
\end{cases}
\)
具有“性质 \(P\)”，求所有满足条件的实数 \(a,b\)；
\item 已知 \(f(x)\) 的值域为 \([0,1)\)，且在 \([0,+\infty)\) 上是严格增函数，证明：\(f(x)\) 是偶函数的充要条件是 \(f(x)\) 具有“性质 \(P\)”.
\end{enumerate}

\end{problem}

\begin{solution}
（1）不具有. 取 \(x_1=-2,x_2=-3\)，则 \(|x_1|<|x_2|\)，但 \(e^{-2}>e^{-3}\)，与性质 \(P\) 矛盾.

（2）由性质 \(P\)，取 \(x_1=0\)，\(x_2>0\)，得 \(b\ge0\)；再结合 \(x_1=0\)，\(x_2<0\) 的比较，可得 \(a<0\). 若 \(b>0\)，可取一组正负对比点推出矛盾，因此 \(b=0\). 此时
\[
f(x)=
\begin{cases}
ax,&x\le0,\\
x,&x>0.
\end{cases}
\]
再比较一正一负两点，最终只能有 \(a=-1\). 故
\[
a=-1,\quad b=0.
\]

（3）充分性：若 \(f\) 是偶函数，则 \(f(x)=f(|x|)\). 又 \(f\) 在 \([0,+\infty)\) 上严格增，所以 \(|x_1|<|x_2|\Rightarrow f(|x_1|)<f(|x_2|)\)，即 \(f(x_1)<f(x_2)\)，故 \(f\) 具有性质 \(P\).

必要性：若 \(f\) 具有性质 \(P\)，设存在 \(x_0>0\) 使 \(f(x_0)\ne f(-x_0)\). 不妨设 \(f(x_0)>f(-x_0)\). 由值域为 \([0,1)\) 且在 \([0,+\infty)\) 上严格增，可取适当 \(x\) 使得 \(|x|<x_0\) 却与性质 \(P\) 矛盾，因此必有 \(f(x_0)=f(-x_0)\). 故 \(f\) 为偶函数.

综上，\(f(x)\) 是偶函数的充要条件是 \(f(x)\) 具有性质 \(P\).
\end{solution}
